
\section{Appendix B. Implementation Details}

\subsection{B.1 Module Structure}

\begin{lstlisting}
@ATTENTION\_REGISTRY.register("infini")
class InfiniAttention(BaseAttention):
    def \_\_init\_\_(self, config, rank=64, compress\_dim=128,
                 local\_window=128, pretrained\_attn=None, **kwargs):
        super().\_\_init\_\_(config, rank)
        
        # Dimensions
        self.dim = 768
        self.num\_heads = 12
        self.head\_dim = 64
        self.all\_head\_size = 768
        self.compress\_dim = 128
        self.local\_window = 128
        
        # Projections with LoRA
        self.q\_proj = LoRALinear(768, 768, lora\_rank=16)
        self.k\_proj = LoRALinear(768, 768, lora\_rank=16)
        self.v\_proj = LoRALinear(768, 768, lora\_rank=16)
        self.out\_proj = LoRALinear(768, 768, lora\_rank=16)
        
        # Compressive memory
        self.memory = nn.Parameter(torch.randn(1, 12, 128, 64) * 0.001)
        self.beta = nn.Parameter(torch.tensor([-5.0]))
        
        if pretrained\_attn is not None:
            self.copy\_weights(pretrained\_attn)
\end{lstlisting}

\subsection{B.2 Forward Pass}

\begin{lstlisting}
def forward(self, hidden\_states, attention\_mask=None, **kwargs):
    B, L, D = hidden\_states.shape
    H, Dh = self.num\_heads, self.head\_dim
    
    # 1. Projections
    q = self.q\_proj(hidden\_states).view(B, L, H, Dh).transpose(1, 2)
    k = self.k\_proj(hidden\_states).view(B, L, self.num\_kv\_heads, Dh).transpose(1, 2)
    v = self.v\_proj(hidden\_states).view(B, L, self.num\_kv\_heads, Dh).transpose(1, 2)
    
    # 2. RoPE
    if position\_embeddings is not None:
        cos, sin = position\_embeddings
        q, k = apply\_rotary\_pos\_emb(q, k, cos, sin)
    
    # 3. GQA
    if self.kv\_repeat \textgreater{} 1:
        k = repeat\_kv(k, self.kv\_repeat)
        v = repeat\_kv(v, self.kv\_repeat)
    
    # 4. Memory attention
    mem = self.memory.expand(B, -1, -1, -1)  # (B, H, d\_c, Dh)
    mem\_scores = torch.matmul(q.float(), mem.float().transpose(-2, -1)) / math.sqrt(Dh)
    mem\_attn = torch.softmax(mem\_scores, dim=-1)
    mem\_out = torch.matmul(mem\_attn, mem.float()).to(v.dtype)
    
    # 5. Local attention
    scale = 1.0 / math.sqrt(Dh)
    scores = torch.matmul(q.float(), k.float().transpose(-2, -1)) * scale
    
    # Sliding window mask
    if self.is\_decoder or self.is\_causal:
        idx = torch.arange(L, device=q.device)
        causal = idx.unsqueeze(1) \textgreater{}= idx.unsqueeze(0)
        window = (idx.unsqueeze(1) - idx.unsqueeze(0)).abs() \textless{} self.local\_window
        window\_mask = causal & window
        scores = scores.masked\_fill(~window\_mask.unsqueeze(0).unsqueeze(0), -1e4)
    
    # Padding mask
    if attention\_mask is not None:
        if attention\_mask.dim() == 2:
            valid\_mask = get\_token\_valid\_mask(attention\_mask, B, L, hidden\_states.device)
            scores = scores.masked\_fill(~valid\_mask.unsqueeze(1).unsqueeze(2), -1e4)
        elif attention\_mask.dim() == 4:
            scores = scores + attention\_mask.clamp(min=-1e4)
    
    attn\_local = F.softmax(scores, dim=-1).to(v.dtype)
    attn\_local = torch.nan\_to\_num(attn\_local, nan=0.0)
    attn\_local = self.dropout(attn\_local)
    
    local\_out = torch.matmul(attn\_local, v)
    
    # 6. Gated fusion
    beta = torch.sigmoid(self.beta)
    combined = local\_out + beta * mem\_out
    combined = combined.transpose(1, 2).contiguous().view(B, L, self.all\_head\_size)
    
    # 7. Output projection
    if self.is\_nested\_encoder:
        return combined, None
    out = self.out\_proj(combined)
    return out, None
\end{lstlisting}

\subsection{B.3 Memory Initialization}

\begin{lstlisting}
self.memory = nn.Parameter(torch.randn(1, H, d\_c, d\_h) * 0.001)
\end{lstlisting}

The scale factor \texttt{0.001} ensures the memory's initial output magnitude is 
small, so even with $\sigma(\beta) \approx 0.007$, the contribution is 
negligible.

\subsection{B.4 Gate Initialization}

\begin{lstlisting}
self.beta = nn.Parameter(torch.tensor([-5.0]))
\end{lstlisting}

At $\beta = -5$: $\sigma(-5) \approx 0.0067$, so memory contributes \textless{}1\% 
initially.

\subsection{B.5 Sliding Window Mask}

For causal models:

\begin{lstlisting}
idx = torch.arange(L, device=q.device)
causal = idx.unsqueeze(1) \textgreater{}= idx.unsqueeze(0)  # j \textless{}= i
window = (idx.unsqueeze(1) - idx.unsqueeze(0)).abs() \textless{} W
mask = causal & window
\end{lstlisting}

For bidirectional models, only \texttt{window} is used.

\subsection{B.6 Common Pitfalls}

\subsubsection{B.6.1 Memory Scale Too Large}

If \texttt{memory} is initialized with \texttt{torch.randn(...)} without scaling, the 
initial memory output is large, and even $\sigma(-5) = 0.007$ multiplied by 
it can disrupt the pretrained behavior.

\subsubsection{B.6.2 $\beta$ Init at 0}

At $\beta = 0$, $\sigma(0) = 0.5$. The memory contributes 50\% at start, 
disrupting the pretrained local attention.

\subsubsection{B.6.3 Missing Window Mask}

Without the window mask, local attention becomes full attention -- the 
whole point is lost.

\subsubsection{B.6.4 Nested Encoder Mis-Set}

For RoBERTa, \texttt{is\_nested\_encoder = True}. Otherwise double projection.
\begin{lstlisting}

---
\end{lstlisting}